\documentclass[../../main]{subfiles}

\begin{document}

\section{スタイルファイルのデモ}

今回分離した \code{.sty} ファイルが正常に動作し\c 各種コマンドが使えることを確認するためのデモです\p 

\subsection{数学スタイル (\code{math\_setup.sty})}
\code{math\_setup.sty} で定義されている独自コマンドや数式機能をテストします\p 
行列表式における行列式 \(\detMat A\) や\c 微小量 \(\del x\) などを用います\p 
また\c レジュメクラスから移行してきたベクトル記法 \(\bvec{v}\) も利用可能です\p 

\begin{equation}
  \Laplace \{ f(t) \} = F(s)
\end{equation}

複素関数における留数や実部・虚部は次のように表記できます\p 
\begin{equation}
  \Residue_{z=c} f(z) \qquad \Rea(z) + i \Ima(z)
\end{equation}

\subsection{単位スタイル (\code{unit\_setup.sty})}
独自定義された物理単位用コマンドのテストです\p 以下のように簡潔に入力が可能です\p 
\begin{itemize}
    \item 長さ: \( 5.0 \, \Met \)
    \item 質量: \( 10 \, \Ki\Gra \)
    \item 力: \( 2.0 \, \New \)
    \item 周波数: \( 60 \, \Hel \)
\end{itemize}

これらの裏側では \code{siunitx} パッケージも読み込まれているため\c そちらの標準的な記法も併用可能です(例: \SI{10}{\meter\per\second})\p 
※注意：新しい \code{siunitx} では \code{\textbackslash qty} 等を推奨していますが\c 互換性のため従来記法も使用可能です\p 

\subsection{TikZスタイル (\code{tikz\_setup.sty})}
多くの \code{usetikzlibrary} があらかじめ読み込まれているため\c 複雑な図もすぐに描画できます\p 以下は影付きノードと矢印のデモです\p 

\begin{figure}[H]
  \centering
  \begin{tikzpicture}[>=stealth, auto, node distance=3cm]
    \node [circle, draw, drop shadow, fill=blue!10] (A) {Node A};
    \node [rectangle, draw, drop shadow, fill=red!10, right of=A] (B) {Node B};
    
    \draw [->, thick] (A) edge[bend left] node {Next} (B);
    \draw [<-, thick, dashed] (A) edge[bend right] node[below] {Back} (B);
  \end{tikzpicture}
  \caption{TikZのライブラリ(shadows等)を用いたテスト描画}
\end{figure}

以上で\c 各スタイルの独立した読み込みと利用方法のデモを終わります\p 
\end{document}
